بروز خطا در این سناریو ناشی از پدیده \lr{Mental Translation} و نبود یک درک یکپارچه از مفاهیم است. وقتی تیم فنی اصطلاح کسب‌وکار (\lr{Customer}) را به دلخواه خود به اصطلاح دیگری (\lr{Client}) در کد نگاشت می‌کند، تفاوت‌های ظریف و حیاتی دامنه از بین می‌رود. از آنجا که طبق منطق کسب‌وکار یک \lr{Customer} (شخصیت حقوقی یا پرداخت‌کننده) می‌تواند شامل چندین \lr{Client} (کاربر نهایی) باشد، اگر توسعه‌دهندگان این دو مفهوم مجزا را در کد قالب یک شیء واحد در نظر بگیرند، قوانین تجاری دچار فروپاشی می‌شوند. به عنوان مثال، فرآیندهایی مانند صدور فاکتور یا بررسی محدودیت‌های مالی که باید در سطح \lr{Customer} اعمال شوند، به اشتباه در سطح تک‌تک \lr{Client}ها ارزیابی شده و منجر به باگ‌های منطقی و محاسباتی مخرب در سیستم می‌شوند.

مفهوم \lr{Ubiquitous Language} این مشکل را با اجبار کردن تمامی اعضا به استفاده از یک واژگان واحد، دقیق و توافق‌شده حل می‌کند. برای برطرف کردن این سوءتفاهم در سطح تیم، \lr{Domain Experts} و توسعه‌دهندگان باید جلسات مشترکی برگزار کنند تا یک واژه‌نامه صریح بسازند و هرگونه ابهام را برطرف کنند. قانون اصلی این است که ترجمه کردن کلمات بین تیم‌ها باید کاملا متوقف شود؛ اگر کسب‌وکار بین پرداخت‌کننده و کاربر نهایی تفاوت قائل است، تیم فنی نیز ملزم است دقیقا از همان کلمات \lr{Customer} و \lr{Client} با همان تعاریف تجاری در تمامی مکالمات روزمره، مستندات و جلسات طراحی استفاده کند.

در سطح مدل دامنه نیز این توافق زبانی باید مستقیما و بی‌واسطه در سورس‌کد پیاده‌سازی شود (\lr{Model-Driven Design}). تیم توسعه باید کدها را \lr{Refactor} کرده و دو مفهوم کاملا مجزا در مدل ایجاد کند. کلاس \lr{Customer} باید مدل‌سازی شود تا صرفا دربرگیرنده منطق مالی، قراردادها و پرداخت‌ها باشد و در کنار آن کلاس \lr{Client} برای مدیریت ویژگی‌های مربوط به دسترسی و استفاده از سیستم طراحی گردد. همچنین ارتباط ساختاری بین این دو (مثلا اینکه یک \lr{Customer} به عنوان یک ریشه تجمیع، مدیریت لیستی از \lr{Client}های زیرمجموعه خود را بر عهده دارد) باید به روشنی در کد پیاده شود تا نرم‌افزار به آینه‌ای تمام‌نما از واقعیت کسب‌وکار تبدیل گردد.